\section{Full problem setup and weighted POD construction}
\label{app:problem_pod}

\subsection{Parameterized full-order model, gauge, and regime coverage}

After spatial discretization, the velocity
$\bm u(t;\mu)\in\mathbb R^{N_u}$ and pressure
$p(t;\mu)\in\mathbb R^{N_p}$ satisfy
\begin{equation}
    \dot{\bm u}
    =
    \bm f_h(\mu)
    +\mathcal L_h(\mu)\bm u
    -\mathcal N_h(\bm u,\bm u;\mu)
    -\nabla_h(\mu)p,
    \qquad
    \operatorname{div}_h(\mu)\bm u=0.
    \label{eq:app_fom}
\end{equation}
Pressure is defined only up to an additive constant.
With positive finite-volume weights $\bm w$, each pressure snapshot is mapped
to the weighted zero-mean gauge
\begin{equation}
    \Gamma_p(q)
    =
    q-\frac{\bm w^\top q}{\bm w^\top\bm 1}\bm 1,
    \qquad
    p^\circ=\Gamma_p(p),
    \qquad
    \bm w^\top p^\circ=0.
    \label{eq:app_pressure_gauge}
\end{equation}

The parameter domain is covered by overlapping regime-local parameter regions,
\begin{equation}
    \mathcal M
    \subseteq
    \bigcup_{r\in\{S,H,P\}}\mathcal M_r,
    \qquad
    \mathcal O_{SH}
    =
    \mathcal M_S\cap\mathcal M_H,
    \qquad
    \mathcal O_{HP}
    =
    \mathcal M_H\cap\mathcal M_P,
    \label{eq:app_regime_cover}
\end{equation}
where $S$, $H$, and $P$ denote the steady, Hopf-onset, and periodic
regimes, respectively.
Only adjacent overlaps are used; no direct $S$--$P$ overlap is considered.


\subsection{Area-weighted local POD}

The finite-volume quadrature induces the weighted inner products
\begin{equation}
    \langle\bm v,\bm z\rangle_{M_u}
    =
    \bm v^\top M_u\bm z,
    \qquad
    \langle q,s\rangle_{M_p}
    =
    q^\top M_p s,
    \qquad
    M_p=\operatorname{diag}(\bm w),
    \qquad
    M_u=\operatorname{diag}(M_p,M_p).
    \label{eq:app_weighted_inner_products}
\end{equation}

For chart $r$, let $d_u^{(r)}$ and $d_p^{(r)}$ denote the retained velocity
and pressure dimensions.
Let $X_u^{(r)}$ and $X_p^{(r)}$ collect the centered velocity and
gauge-corrected pressure snapshots from the training trajectories assigned
to that chart.
The corresponding chart-specific means are denoted by
$\bar{\bm u}^{(r)}$ and $\bar p^{(r),\circ}$.
The local POD bases solve
\begin{align}
    \Phi_u^{(r)}
    &\in
    \underset{\Psi^\top M_u\Psi=I_{d_u^{(r)}}}
    {\operatorname{arg\,min}}
    \left\|
        X_u^{(r)}
        -\Psi\Psi^\top M_uX_u^{(r)}
    \right\|_{M_u,F}^2,
    \\
    \Phi_p^{(r)}
    &\in
    \underset{\Psi^\top M_p\Psi=I_{d_p^{(r)}}}
    {\operatorname{arg\,min}}
    \left\|
        X_p^{(r)}
        -\Psi\Psi^\top M_pX_p^{(r)}
    \right\|_{M_p,F}^2,
    \label{eq:app_pod_optimization}
\end{align}
where
\begin{equation}
    \|Y\|_{M,F}^2=\operatorname{tr}(Y^\top MY).
\end{equation}

Equivalently, consider the weighted singular value decompositions
\begin{equation}
    M_u^{1/2}X_u^{(r)}
    =
    U_u^{(r)}\Sigma_u^{(r)}{V_u^{(r)}}^\top,
    \qquad
    M_p^{1/2}X_p^{(r)}
    =
    U_p^{(r)}\Sigma_p^{(r)}{V_p^{(r)}}^\top.
\end{equation}
The retained bases are
\begin{equation}
    \Phi_u^{(r)}
    =
    M_u^{-1/2}
    U_u^{(r)}(:,1{:}d_u^{(r)}),
    \qquad
    \Phi_p^{(r)}
    =
    M_p^{-1/2}
    U_p^{(r)}(:,1{:}d_p^{(r)}),
    \label{eq:app_weighted_svd_bases}
\end{equation}
and satisfy
\begin{equation}
    {\Phi_u^{(r)}}^\top M_u\Phi_u^{(r)}=I,
    \qquad
    {\Phi_p^{(r)}}^\top M_p\Phi_p^{(r)}=I.
\end{equation}
Because the pressure snapshots and their chart mean satisfy the same weighted
gauge condition, the pressure basis also satisfies
\begin{equation}
    \bm w^\top\Phi_p^{(r)}=\bm 0^\top .
\end{equation}

Encoding and decoding are performed independently within each chart:
\begin{align}
    \bm a^{(r)}
    &=
    {\Phi_u^{(r)}}^\top
    M_u
    \bigl(\bm u-\bar{\bm u}^{(r)}\bigr),
    &
    \mathcal D_r^u(\bm a^{(r)})
    &=
    \bar{\bm u}^{(r)}
    +\Phi_u^{(r)}\bm a^{(r)},
    \\
    \bm b^{(r)}
    &=
    {\Phi_p^{(r)}}^\top
    M_p
    \bigl(\Gamma_p(p)-\bar p^{(r),\circ}\bigr),
    &
    \mathcal D_r^p(\bm b^{(r)})
    &=
    \bar p^{(r),\circ}
    +\Phi_p^{(r)}\bm b^{(r)}.
    \label{eq:app_encode_decode}
\end{align}

The reduced coefficients are standardized using chart-local statistics
estimated from the corresponding training trajectories:
\begin{equation}
    \widetilde{\bm a}^{(r)}
    =
    \bigl(
        \bm a^{(r)}-\bm\mu_a^{(r)}
    \bigr)
    \oslash
    \bm\sigma_a^{(r)},
    \qquad
    \widetilde{\bm b}^{(r)}
    =
    \bigl(
        \bm b^{(r)}-\bm\mu_b^{(r)}
    \bigr)
    \oslash
    \bm\sigma_b^{(r)}.
    \label{eq:app_local_scalers}
\end{equation}
Additional specialist features and their normalization are specified in
Appendix~\ref{app:specialists}.

All chart-local means, POD bases, and scaling statistics are constructed
independently.
Consequently, reduced coefficient vectors from different charts are not
assumed to be componentwise comparable; only their decoded physical fields lie in a common physical comparison space.






\section{Projected Galerkin and pressure operators}
\label{app:rom}

\subsection{Chart-local Galerkin momentum dynamics}

For each chart $r\in\{S,H,P\}$, substituting the chart-local affine reconstructions
from Appendix~\ref{app:problem_pod} into the semi-discrete
momentum equation and testing against the local velocity basis yields
\begin{equation}
    \dot{\bm a}^{(r)}
    =
    \bm c_r(\mu)
    +
    A_r(\mu)\bm a^{(r)}
    +
    H_r(\mu):
    \bigl(
        \bm a^{(r)}\otimes\bm a^{(r)}
    \bigr)
    +
    C_r^p(\mu)\bm b^{(r)}.
    \label{eq:app_projected_momentum}
\end{equation}
Here ``$:$'' denotes the double contraction.
For velocity modes
$\bm\phi_i^{u,(r)}$,
$i=1,\ldots,d_u^{(r)}$,
and pressure modes
$\phi_\ell^{p,(r)}$,
$\ell=1,\ldots,d_p^{(r)}$,
the projected operators are
\begin{align}
    [\bm c_r(\mu)]_i
    &=
    \Bigl\langle
        \bm\phi_i^{u,(r)},
        \bm f_h(\mu)
        +
        \mathcal L_h(\mu)\bar{\bm u}^{(r)}
        -
        \mathcal N_h
        \bigl(
            \bar{\bm u}^{(r)},
            \bar{\bm u}^{(r)};
            \mu
        \bigr)
        -
        \nabla_h(\mu)\bar p^{(r),\circ}
    \Bigr\rangle_{M_u},
    \\
    [A_r(\mu)]_{ij}
    &=
    \Bigl\langle
        \bm\phi_i^{u,(r)},
        \mathcal L_h(\mu)\bm\phi_j^{u,(r)}
        -
        \mathcal N_h
        \bigl(
            \bar{\bm u}^{(r)},
            \bm\phi_j^{u,(r)};
            \mu
        \bigr)
        -
        \mathcal N_h
        \bigl(
            \bm\phi_j^{u,(r)},
            \bar{\bm u}^{(r)};
            \mu
        \bigr)
    \Bigr\rangle_{M_u},
    \\
    [H_r(\mu)]_{ijk}
    &=
    -
    \Bigl\langle
        \bm\phi_i^{u,(r)},
        \mathcal N_h
        \bigl(
            \bm\phi_j^{u,(r)},
            \bm\phi_k^{u,(r)};
            \mu
        \bigr)
    \Bigr\rangle_{M_u},
    \\
    [C_r^p(\mu)]_{i\ell}
    &=
    -
    \Bigl\langle
        \bm\phi_i^{u,(r)},
        \nabla_h(\mu)\phi_\ell^{p,(r)}
    \Bigr\rangle_{M_u}.
    \label{eq:app_projected_operators}
\end{align}

We denote the resulting chart-local projected momentum vector field by
\begin{equation}
    \mathcal F_r^G
    \bigl(
        \bm a^{(r)},
        \bm b^{(r)};
        \mu
    \bigr)
    =
    \bm c_r(\mu)
    +
    A_r(\mu)\bm a^{(r)}
    +
    H_r(\mu):
    \bigl(
        \bm a^{(r)}\otimes\bm a^{(r)}
    \bigr)
    +
    C_r^p(\mu)\bm b^{(r)}.
    \label{eq:app_galerkin_backbone}
\end{equation}
This vector field provides the explicit projected Galerkin backbone of the
regime-local specialist introduced in
Eq.~(\plaineqref{eq:main_specialist}).


\subsection{Gauge-consistent pressure--Poisson map}

The pressure component is represented through a chart-local
Pressure--Poisson relation.
At the continuous level, the gauge-corrected pressure satisfies
\begin{equation}
    -\Delta p^\circ
    =
    \nabla\!\cdot
    \left[
        (\bm u\!\cdot\!\nabla)\bm u
        -
        \bm f
    \right],
    \qquad
    \int_\Omega p^\circ\,\mathrm d\Omega=0,
    \label{eq:app_continuous_pressure_poisson}
\end{equation}
with boundary conditions inherited from the momentum equation.
After discretization, substitution of the chart-local affine velocity
representation, and projection onto the local pressure basis, the reduced
pressure coefficients satisfy
\begin{equation}
    K_r^p(\mu)\bm b^{PP,(r)}
    =
    \bm d_r(\mu)
    +
    B_r(\mu)\bm a^{(r)}
    +
    Q_r(\mu):
    \bigl(
        \bm a^{(r)}\otimes\bm a^{(r)}
    \bigr).
    \label{eq:app_pressure_poisson_system}
\end{equation}
Provided that the gauge-restricted reduced system is nonsingular, the corresponding chart-local Pressure--Poisson map is
\begin{equation}
    \mathcal Q_r^{PP}
    \bigl(
        \bm a^{(r)};
        \mu
    \bigr)
    =
    [K_r^p(\mu)]^{-1}
    \left[
        \bm d_r(\mu)
        +
        B_r(\mu)\bm a^{(r)}
        +
        Q_r(\mu):
        \bigl(
            \bm a^{(r)}\otimes\bm a^{(r)}
        \bigr)
    \right].
    \label{eq:app_pressure_poisson_map}
\end{equation}

Because the pressure basis lies in the weighted zero-mean subspace defined in
Appendix~\ref{app:problem_pod}, the constant-pressure null direction is
removed before projection. Removing that null direction alone does not establish nonsingularity of an arbitrary discretized reduced operator. The inverse notation above denotes the reduced linear solve, not an instruction to form an explicit inverse. The linear term includes cross interactions with the nonzero chart mean; boundary contributions depend on the chosen discretization. This projected relation describes the operator construction and should not be interpreted as an exact nonlinear constraint on the learned pressure correction or on fused outputs.
Thus $\mathcal Q_r^{PP}$ produces pressure coefficients in the native
coordinates of chart $r$ without introducing an additional cross-chart
pressure gauge.

In the archived circular-cylinder Periodic runtime, the pressure base is evaluated from preassembled constant, linear, and quadratic coefficient arrays, rather than by solving a full-mesh pressure equation at each prediction step. This runtime polynomial evaluation alone does not establish how its offline coefficient arrays were obtained; intrusive projection and data fitting must be distinguished when interpreting the operator provenance.

The pair
\begin{equation}
    \left(
        \mathcal F_r^G,
        \mathcal Q_r^{PP}
    \right)
\end{equation}
is constructed independently for each regime-local chart.
Accordingly, each chart carries its own parameter region, mean fields,
POD bases, coefficient scalings, projected momentum operators, and
Pressure--Poisson map.
These chart-local operators are used only within their native reduced
coordinates; no componentwise interpolation of reduced states or projected
operators is assumed across charts.
